\documentclass[11pt]{article}
\usepackage[T1]{fontenc}
\usepackage{lmodern,microtype}
\usepackage{amsmath,amssymb,booktabs,enumitem}
\usepackage[letterpaper,margin=1in]{geometry}
\usepackage[colorlinks=true,urlcolor=blue,citecolor=blue]{hyperref}
\newcommand{\Q}{\mathbb Q}
\newcommand{\F}{\mathbb F}
\title{Metaplectic Jacquet--Langlands Transfer, Companion Paper (Rank Two, Wild Place)\\Revision Log: Draft 3 to Draft 4}
\author{Prepared by Claude (Anthropic) for David Joyner}
\date{November 8, 2026}
\begin{document}
\maketitle
\section*{Scope}
Draft 4 makes exactly the six corrections requested in ChatGPT's referee
report on Draft 3, with no change of mathematical scope.  All six were
checked and found correct before being applied.  Sections other than
8.2 (one lemma statement), 12.1, 13 are unchanged.
\section*{The six corrections}
\begin{enumerate}[leftmargin=*]
\item \textbf{Central character in the two-group step} (Theorem
\texttt{thm:reduction-all}).  Draft 3 twisted $\rho^\vee$ by
$\theta_1^2\circ\mathrm{Nrd}$ and called the result trivial on $F^\times$
(even case) or of central character $\mathrm{sgn}$ (odd case).  Since
$\mathrm{Nrd}(z)=z^2$, that twist restricts to $\theta_1$ on scalars, so
both claims were wrong.  Following the referee, the proof is now
organized around the genuine order-three character $\psi_P$ of
$\Gamma_{P_0}$ (trivial on a Sylow $2$-subgroup, $\psi_P(z,0)=\theta_1(z)$,
hence $\psi_P(x_\eta)=\theta_1(\eta^4)$) and the untwisted $\rho^\vee$:
\[
 \phi_\rho|_{\Gamma_{P_0}}=\psi_P\otimes\Bigl(3\rho^\vee|_{D_8}+3\dim\rho\sum_{\pi}(\dim\pi)\Theta_\pi\Bigr)
 \quad\text{resp.}\quad
 \psi_P\otimes\bigl(3f\,\Theta_{\rho^\vee|_{SD_{16}}}+3\dim\rho\,(2\chi_1+2\chi_2)\bigr).
\]
The $SD_{16}$ identification and the swap $f\chi_1=\chi_2$ are unchanged
(the referee verified them).
\item \textbf{Brauer case split} (same theorem).  When $T\supset\mu_2$ a
noncyclic $B_2$ can have cyclic image in $D_8$ (e.g.\ $\langle a^4,b\rangle$),
and $C_{\Gamma_0}(B_2)$ then contains noncentral elements of the
$j$-torus, so Draft 3's ``noncyclic $B_2$'' case was wrong.  The split
is now by the image $\bar B_2\subset D_8$: trivial (hypothesis on
$\Phi_\rho$), nontrivial cyclic (toral descent, using
$B_2\subset\langle y\rangle(B_2\cap T)$ and $T\subset T_{E_y}$), noncyclic
(centralizer lemma, then the two-group computation).  Lemma
\texttt{lem:centralizers}(b) is restated for $2$-subgroups with
noncyclic image, which is what its proof always used; for $\rho=\mathbf 1$
($T=\mu_3$) the two formulations agree, so Theorem
\texttt{thm:reduction} is unaffected.
\item \textbf{General-$p$ character formula} (Theorem \texttt{thm:H1-gen}).
The value of $\psi_{2d}$ at the star lift is $\theta_0(z)$ with
$z\equiv\mathrm{Nrd}(X)^{(p+1)/2}$, so the formula is
$\zeta^{e-\kappa(X)}\theta_0(\mathrm{Nrd}X)^{(p+1)/2}$, with $\theta_0^p=1$
now stated; Draft 3's exponent $-1$ is correct only for $p=3$.  The
multiplicativity proof is unchanged.
\item \textbf{Depth lemma} (Lemma \texttt{lem:depth-gen}(c)).  Both cases
are now written out: ramified via the isomorphism
$U_F^{(d)}/U_F^{(d+1)}\to U_E^{(2d)}/U_E^{(2d+1)}$ and bijectivity of the
$p$-th power map $U_E^{(k_0+1)}\to U_E^{(2d+1)}$; unramified via the
boundary map $\bar a\mapsto\bar a_0^{\,p}\wp(\bar a/\bar a_0)$ (with
$p/\varpi^e\equiv-\bar a_0^{\,p-1}$ from $\zeta_p^p=1$), whose image is
$\bar a_0^{\,p}\ker\mathrm{Tr}_{k_E/\F_p}$ and misses $k_F$ because
$\mathrm{Tr}_{k_E/\F_p}=2\,\mathrm{Tr}_{k_F/\F_p}$ on $k_F$ and $p$ is odd.
\item \textbf{Typo} in Theorem \texttt{thm:levelone}: $w^4=\mathrm{Nm}(w)$,
not $\mathrm{Nm}(w)^2$.
\item \textbf{Citation and summation by parts.}  Lemma \texttt{lem:model}(a)
now cites Kubota's 1967 paper for $\mathrm{SL}_2$ and identifies the
$\mathrm{GL}_2$ formula as the $r=2$, $c=0$ Kazhdan--Patterson cocycle
with the $n$-th power symbol (added \texttt{[Kub67]} to the
bibliography).  The summation by parts in Theorem \texttt{thm:pthpower}
is written out: $c_{j-1}-2c_j+c_{j+1}=-1$, $c_1=c_p=0$,
$c_0=c_{p+1}=-p/2\equiv0$, and the four boundary terms are displayed.
\end{enumerate}
\section*{Status}
Theorem \texttt{thm:reduction-all} [P] (repaired); Theorem
\texttt{thm:H1-gen} [P] (formula corrected, lemma completed);
Theorem \texttt{thm:pthpower} [P] (accepted by the referee); Theorem
\texttt{thm:levelone} [P$^c$] (conditional on Target A, as before).
Nothing else changed status.
\end{document}
